\chapter{Kinetic Theory: Classical, Relativistic, and Quantum}
\section{Math}
\subsection{Eigenfunction expansion and resolvent}
An example of the application of the Liouville equation integrating factor is an ideal gas.

Taking the purely kinetic Hamiltonian

\[
H = \sum_{s=1}^N \frac{p^2_s}{2m}, \quad 0 \leq x^{i}_s \leq L
  \]
  Where $ x^{i}_s$ denotes the  Cartesian component of $\mathbf{x}_s$. Let us call the $\vu{\Lambda}$ operator. Thus:

\[
\vu{\Lambda}_0 = -i \sum_s \pdv{H}{\mathbf{p}_s} \cdot \pdv{}{\mathbf{x}_s} = -i \sum_s \mathbf{v}_s \cdot \pdv{}{\mathbf{x}_s}
  \]

  The Eigenfunctions of $\vu{\Lambda}$ satisfy the eigenvalue equation.

Taking this and solving we get
\[
D(\mathbf{x}^n, \mathbf{p}^N, t) = \frac{1}{L^{3N/2}}\sum_k D_k \mathbf{p}^Ne^{i \sum_s k_s \cdot (\mathbf{x}_s-\mathbf{v}t)}
  \]

  Taking a free particle propagator we find that one evolves the function through time via
\[
  \exp{-t \mathbf{v_s}} \cdot \pdv{}{\mathbf{x}_s} \]

One last solution technique presented to initial value problems for the Liouville equation uses the Laplance transform operated upon the Liouville equation
\[
\int_0^\infty dte^{-st}\left(i \pdv{D}{t} = \vu{\Lambda}D  \right)
  \]

  From there we define
\[
\vu{R} = (\vu{\Lambda} - is)^{-1}
  \]
  as the resolvent operator.

\subsection{Distribution Functions}

Nothing

\subsection{Markov Process}

Given the phase vector $z = (1,2,...,N)$, in terms of this variable, the two-time, or conditional, distribution function may be written at the product
\[
\prod (z,t \mid z_0, t_0) dz
  \]

  This distribution has three fundemental principles
\begin{itemize}
  \item The system undergoes a transition from some state in the interval $t-t_0$
  \item The system does not change its state in zero time
  \item The function $\Pi$ obeys the chapman-kolmogorov equation
\end{itemize}
\[
  \Pi (z,t \mid z_0, t_0) = \int  \Pi (z,t \mid z', t') \Pi (z',t' \mid z_0, t_0)dz'
\]
The integral represents a sum over intermediate states.

For systems whose behavior is homogeneous in time, processes can only depend on the time interval and not the initial or final times.

Our aim is then to obtain a differential equation for $\Pi$ appropriate to the case where phase variables go over a countable set.

The integral equation becomes
\[
\Pi( l \mid l_0 ; t+ \delta t ) = \sum_{\forall j} \Pi ( l \mid j ; \delta t ) \Pi ( j \mid l_0 ; t )
  \]

  Conditions 1 and two become
\[
\sum_{\forall l}\Pi (l \mid l_0; t) =1
  \]
\[
\Pi(l \mid l_0;0) = \delta_{ll_0}
  \]
  For $l \neq j$
\[
\Pi (l \mid j; \delta t) \to 0 \, \text{as} \, \delta t \to 0
  \]

  With some further derivation I won't go into we get the canonical form of the mater equation
\[
  \pdv{\Pi(l \mid l_0 ; t)}{t} = \sum_j (w_{jl} \Pi (j \mid l_0 ; t) - w_{lj} \Pi (l \mid l_0; t))
\]

Which is essentially this psuedo equation
$$\text{Change} = \text{Gain} - \text{Loss}$$

\subsection{Central Limit Theorem}

Variance is writen as
\[
D(\xi) = \sum_{\forall x}(x - \langle x \rangle)^2 P(x)
  \]

  The characteristic function being
\[
  \phi(a) = \sum_{\forall x} P(x)e^{iax}  \]

From here we can take a Fourier transformation expand into a series that will be used later. We can then define a random walk.

An importaint distribution is defined as
\[
P(l,n) = \frac{n !}{r!(n-r)!}\mathfrak{p}^r(1- \mathfrak{p})^{n-r}
  \]
  When we define $r \equiv \frac{n+l}{2}$.

This lets us use some math to define
\[
    P(r) \equiv = P(r, \frac{\lambda}{n} = \frac{\lambda '}{r!}e^{-\lambda}  \]

More commonly known as the Poisson distribution.

\section{Analyses of the Liouville Equation}

\subsection{BBKGY Hierarchy}

Generalizing what I have already done.

\subsection{Correlation Extrapolation: Vlasov limit}
If we assume $N \gets s$, then we can rewrite as
\[
\left(\pdv{}{t} - \vu{L_s}  \right)f_s = \vu{I}_s f_{s+1}
  \]

  With some further math we can define both how dense the fluid is and how strongly coupled.

  Vlasov equation results in the a well defined limit, that can be caused by plasma comprised of electrons in a neutralizing, fixed, positive background.

\subsection{Diagrams: Prigogine Analysis}
Ilya Prigogine's work on nonequilibrium thermodynamics
investigates how irreversible processes generate entropy
and how systems maintained away from thermodynamic
equilibrium can develop organized structures.

The central concepts include:
\begin{itemize}
    \item Entropy production and irreversible processes.
    \item Thermodynamic equilibrium versus nonequilibrium.
    \item Linear irreversible thermodynamics.
    \item Minimum entropy production under restricted conditions.
    \item Instabilities and dissipative structures.
    \item The emergence of macroscopic organization from
          microscopic dynamics.
\end{itemize}

A key distinction is between a system's total entropy
and its internal entropy production. An open system can
maintain or develop organization while exporting entropy
to its environment.

A typical workflow
\begin{enumerate}
    \item Define the system, its boundaries, and its
          exchanges of matter and energy.
    \item Identify state variables and conserved quantities.
    \item Derive the entropy balance and identify the
          internal entropy production.
    \item Specify thermodynamic forces and fluxes.
    \item Establish constitutive relations and check
          thermodynamic consistency.
    \item Determine stationary solutions under the
          imposed boundary conditions.
    \item Analyze stability as control parameters vary.
    \item Identify possible bifurcations, patterns,
          oscillations, or other nonequilibrium states.
    \item Determine whether the system lies within the
          regime of validity of linear irreversible
          thermodynamics.
\end{enumerate}


The central insight is that irreversible dynamics need
not produce only featureless relaxation. Under suitable
conditions, sustained energy or matter fluxes can support
organized nonequilibrium states. Their existence and
stability must be established from the governing
dynamics rather than inferred from entropy production
alone.

\subsection{Bogoliubov Hypothesis}

For context: In the interval $\tau_1$, two molecules are in each others interaction domain. The interval $\tau_2$ is the mean-free-collision times, which is the mean times between collisions. The time $tau_3$ is the average time taken for a molecule to traverse the container in which the gas is confined.

This section for the most part finds the thermodynamic limit and fluid limit.

They also derive the Boltzmann equation.

So these Bogoliubov and Prigogin techniques give us the ability to create kinetic equations at higher configurations.

\subsection{Klimontovich Picture}

In this section we turn to an alternative description of kinetic theory formulated by Klimontovich. In this formalism, phase densities are introduced, moments of which are found to be related to multiparticle distribution functions introduced earlier. Moments of deviations from the mean of these phase densities are similarly found to be related to previously encountered correlation functions. Applying these algorithms to an equation of motion for the first moment of the phase densities is found the first equation of the hierarchy.

\subsection{Grad's Analysis}

Starting with Liouville equation written in this form
\[
\pdv{f_N}{t} + \sum^N \pdv{}{x_i} \cdot (v_i f_N) + \frac{1}{m} \sum^N G_i \cdot \pdv{f_N}{v_i} = 0
  \]

  \section{The Boltzmann Equation, Fluid
Dynamics, and Irreversibility}

\subsection{Fluid Dynamic Equations and the Boltzmann H
Theorem}

\subsection{Transport Coefficients}


Consider that a fluid is in equilibrium (or steady state) at given uniform values of
n, u, and T . Gradients of these variables develop under perturbation. The fluid
responds to these gradients in a manner to restore equilibrium (Le Chatelier’s
principle). Thus, with R written for “response” we may write

\[
R[\nabla n] = nu
  \]
\[
R[\nabla u] = S
  \]
\[
R[\nabla T] = Q
  \]
\[
R[\nabla \Phi] = J
  \]
  That is, fluid motion will develop in response to gradients in density.

  The coefficients that relate gradients of the perturbation to response motions are called transport coefficients.

\subsection{The Chapman–Enskog Expansion}

The collision integral of the Boltzmann equation
\[
J(F \vert F) \equiv -F \int \int \sigma d \Omega g dv_1 F(v_1) + \int \int \sigma d \Omega g dv_1 F' F_1'
  \]

  The quantity
\[
\nu = \int \int \sigma d \Omega g dv_1 F(v_1)
  \]

  Represents collision frequency

  Allows us to rewrite the Boltzmann Equation


  Which serves to define the collision-integral operator $vu{I}$. Which has a dimensions of frequency.

  We can see that a large collision frequency is equivalent to a small mean free path.

  We can once again rewrite the Boltzmann equation as
\bq
\left(\pdv{}{t} + \vu{D}  \right)F = \frac{1}{\epsilon} \vu{I}F
\label{Bolt_Chap}
\eq

Step 2 is to introduce the expansion
\bq
F = F^{(0)} + \epsilon F^{(2)} + \epsilon^2 F^{(2)} + ...
\label{series_F}
\eq

The normalization of F satisfies the moment relations
\bq
n = \int F dv, \quad n \vb{u} = \int F \vb{v} d\vb{v}, \quad \frac 3 2 n k_B T = \int \frac{mc^2}{2}F d\vb{v}
\eq

Step three stipulates that the variables are all $0(1)$ quantities and stem from $F^{(0)}$. Where the rest of the series \eqref{series_F} correspond to the higher moments, I won't go through and write all of them out.

Using the same method we can expand $\vu{D}$.

We can also expand $\vu{J}$ y
\bq
\vu{hat}(F \vert F) = \vu{J}^{(0)}(F^{(0)}) + \epsilon \vu{J}^{(1)}(F^{(0)}F^{(1)}) + ...
\eq

The fourth steps is to address the time derivatives

The physical meaning of this expansion is that lowest-order terms vary most rapidly, whereas higher-order terms are more slowly varying.

The book then goes on expressing the equation and deriving results, but I think what I have done so far is sufficient.


\subsection{The Druyvesteyn Distribution}
